
@@M0@@
@@M1@@
définissant sur le foncteur @@M2@@ une multiplication et des opérations de @@M3@@. D'après 1.13, la multiplication est définie par sa valeur sur l'élément @@M4@@ de @@M5@@~; elle sera notée @@M6@@~; on a donc :
@@M7@@
Le @@M8@@-ième polynôme universel de cette loi est la composante isobare de poids @@M9@@ de @@M10@@ (1.15) et sera noté @@M11@@.

Les opérations de @@M12@@ seront encore notées par l'exponentielle (3.8). Les polynômes universels définissant ces opérations sont les composantes isobares de plus bas degré des @@M13@@ (1.15)~; ceux-ci étant isobares, ce sont les @@M14@@ eux-mêmes.

Soit @@M15@@~; l'anneau @@M16@@ défini par les lois précédentes est en général sans élément unité~; on étend donc cet anneau par adjonction des scalaires de @@M17@@, en introduisant le nouvel anneau @@M18@@, muni des lois de composition définies par : @@M19@@, @@M20@@,
@@M21@@
@@M22@@

L'anneau ainsi construit admet (1,1) pour l'élément unité, et @@M23@@ s'identifie à l'idéal @@M24@@. Il est muni de façon naturelle d'une structure de @@M25@@-algèbre, telle que :
@@M26@@

